%%%%%%%%%%%%%%%%%%%%%%%%%%%%%%%%%%%%%%%%%%%%%%%%%%%%%%%%%%%%%%%%%%%%%%%%
% Plantilla TFG/TFM
% Escuela Politécnica Superior de la Universidad de Alicante
% Realizado por: Jose Manuel Requena Plens
% Contacto: info@jmrplens.com / Telegram:@jmrplens
%%%%%%%%%%%%%%%%%%%%%%%%%%%%%%%%%%%%%%%%%%%%%%%%%%%%%%%%%%%%%%%%%%%%%%%%

\chapter{Metodología}
\label{metodologia}

\section{Diseño de los experimentos}

/**/

El entrenamiento tanto en los modelos generativos (RNN, LSTM, decoder-only) como en los clasificadores (encoder-only) es business as usual 
(pero explícalo lol)

Para evaluación, en los \textbf{modelos generativos}:

Nuestro caracter de inicio es el \^. Dado ese caracter, los modelos tienen que generar una cadena que lo continúe.

Fácilmente podríamos simplemente comparar la cadena que genera con la cadena que queríamos que generara, pero esto introduce dificultades
bastante únicas a nuestro objeto de estudio.
Por ejemplo, si esperamos $a^4b^4$, y al modelo habiendo aprendido correctamente la gramática, le da por generar $a^7b^7$, 
entonces estaríamos castigándolo por generar una cadena perfectamente válida.

Adicionalmente, también perderíamos bastante nuance relacionada con la correctitud de cada caracter en cada posición de la cadena.
Para un lenguaje $a^Nb^N$, lo que viene después de cualquier ``a" debería ser alguna distribución de probabilidad entre ``a" y ``b" (con suerte
dándole algo de favoritismo a ``a"). En cambio, una vez se genera una ``b", la distribución de probabilidad debería colapsar en casi completo 
favor a ``b", porque la gramática no permite después de una ``b" cualquier cosa que no sea otra ``b", o el fin de cadena.

Entonces, si simplemente decimos si la cadena generada pertenece o no al lenguaje, estamos perdiendo datos sobre lo cercana a correcta que 
era la cadena generada. Si esperamos ``aaabbb", es peor generar ``aaaba" que generar ``aaabb", porque la primera es una fumada y la segunda es
un pequeño desajuste de probabilidades.

De modo que tocará evaluar paso por paso e ignorar la idea de una cadena esperada. Para cada cadena de entrada, los modelos generarán el siguiente
token, y nuestra única manera de evaluarlo será comprobando si ese caracter podía o no ponerse ahí según la gramática. Para esto será necesario
implementar un méŧodo específico para todos los lenguajes, pero vivir es sufrir. Esa será nuestra forma de evaluar el parámetro de accuracy.

En los \textbf{clasificadores}:

Simplemente dada una cadena generarán un flotante entre 0 y 1, con 0 siendo absolutamente NO y 1 siendo absolutamente SÍ. A estos sí que se les 
dará la cadena entera y se esperará un valor de vuelta. Sin más.


